\documentclass[a4paper]{article}
%
%
\addtolength{\textwidth}{3cm}
\addtolength{\oddsidemargin}{-1.5cm}
%
%
\usepackage{flafter}            % Floats appear only after they are referenced
%
% Befehl, um die Caption bei Tabellen abzusetzen
%
\newcommand{\mytabspace}{\medskip}
%
%
\begin{document}
%
%
\section{Some basic notions on XML}
%
%
\section{Meta structure of the parameter file}
%
%
The basic structure of a parameter file for CFS++ in our XML syntax is as
follows:
%
%
\begin{verbatim}
<?xml version="1.0"?>
<cfsSimulation xmlns="http://www.cfs++.org">

 ... Parameters ...

</cfsSimulation>
\end{verbatim}
%
%
The first line gives version information for the XML parser, while the second
line is the starting tag for the main enclosing element called cfsSimulation.
The last line then is the closing tag for this main element. All parameters
must appear between the opening and the closing tag.

As noted in the previous section, the XML parameter file is built up in a
hierarchical fashion. It can be imagined as a book or tree. The cfsSimulation
element is the root of this tree and forms the zeros level. All leaves of this
root are either parameters or the starting nodes of subtrees aka sections or
subsections. Table \ref{tab:Layer1Elements} shows all leaves of the first
level of the data tree. It also gives the type of the individual leaves. Here
plain denotes an entry that represents a certain parameter, while section
indicates that the entry itself is a (sub)section starting a new subtree.

Table \ref{tab:Layer1Elements} also shows which entries are required and which
are optional. Note, however, that the entries, if they appear, must appear
precisely in the order specified in table \ref{tab:Layer1Elements}.
%
%
\begin{table}
\caption{\label{tab:Layer1Elements}Entries of the parameter file on level 1.}
\mytabspace
\centering
\begin{tabular}{lll}
\hline\noalign{\smallskip}
leave name & leave type & occurence constraints \\
\noalign{\smallskip}\hline\noalign{\smallskip}
analysis     & plain   & required\\
geometry     & plain   & required\\
input        & plain   & optional\\
output       & plain   & optional\\
materialData & plain   & optional\\
domain       & section & required \\
pdeList      & section & required \\
transient    & section & required for transient analysis \\
\noalign{\smallskip}\hline
\end{tabular}
\end{table}
%
%
\section{The 'domain' Section}
%
%
\begin{verbatim}
<domain>
  <region name="cantilev" material="silizium"/>
  <nodes name="ux-dof"/>
  <nodes name="uy-dof"/>
  <nodes name="uy-load"/>
  <elements name="dummy"/>
</domain>
\end{verbatim}
%
%
\section{The 'pdeList' Section}
%
%
\section{The 'transient' Section}
%
%
The parameter file must contain a 'transient' section, if the type of analysis
is transient. This section contains information on \ldots
%
%
\begin{verbatim}
<transient>
  <numsteps>    20    </numsteps>
  <firstdt>     5e-05 </firstdt>
  <stepsavebeg> 1     </stepsavebeg>
  <stepsaveend> 20    </stepsaveend>
  <stepsaveinc> 1     </stepsaveinc>
  <timeDataFile name="sin1khz.dat"/>
  <timeDataFile name="sin100khz.dat"/>
</transient>
\end{verbatim}
%
%
\begin{table}
\caption{\label{tab:TransientSection}Entries of 'transient' section.}
\mytabspace
\centering
\begin{tabular}{lll}
\hline\noalign{\smallskip}
entry name & value type & number of occurences \\
\noalign{\smallskip}\hline\noalign{\smallskip}
numsteps     & non-negative integer    & once \\
firstdt      & positive floating point & once \\
stepsavebeg  & non-negative integer    & once \\
stepsaveend  & non-negative integer    & once \\
stepsaveinc  & non-negative integer    & once \\
timeDataFile & string                  & arbitrary \\
\noalign{\smallskip}\hline
\end{tabular}
\end{table}

\end{document}
